%% kapa: exam=1 round=23 year=2012 period=2 form=A subject=appraisal_law subject_ko=부동산관계법규 seq=32 number=72
%% source: Q-Net 감정평가사 기출문제 (cst00309) artlSeq=5052580 file=raw/1008/5052580/2012년 제23회 감정평가사 제1차시험 기출문제(2교시 A형).pdf
%% conversion: kordoc 4.24.51 -> tools/gen1.py (auto); review=proofread-sonnet; flags=none
\begin{question}{72}
건축법령상 건축물의 대지와 도로에 관한 설명으로 옳지 않은 것은?
\begin{choices}
  \item 대지의 배수에 지장이 없는 경우에는 대지가 인접한 도로면보다 낮아도 된다.
  \item 녹지지역의 건축물은 면적이 200\,m$^2$ 이상인 대지에 건축하는 경우에도 조경 등의 조치를 하지 아니할 수 있다.
  \item 건축물의 대지는 2 m 이상이 자동차전용도로가 아닌 도로에 접하여야 하지만, 해당 건축물의 출입에 지장이 없다고 인정되는 경우에는 그러하지 아니하다.
  \item 이해관계인이 해외에 거주하여 동의를 받기 곤란한 경우에 허가권자는 건축위원회의 심의를 거쳐 이해관계인의 동의 없이 도로의 위치를 지정ㆍ공고할 수 있다.
  \item 건축물은 지표의 위ㆍ아래 부분에서 건축선의 수직면을 넘어서는 아니 되며, 도로면으로부터 높이 4.5 m 이하에 있는 창문은 열고 닫을 때 건축선의 수직면을 넘지 않는 구조로 하여야 한다.
\end{choices}
\end{question}
\answer{5}
